\subsection{Little Red Dots as Early Saturation Signatures}
  \label{subsec:lrd_saturation}

  Recent JWST observations have revealed a population of compact,
  luminous, high-redshift sources often referred to as ``Little Red Dots''.
  These objects appear overmassive relative to expectations from standard
  hierarchical assembly scenarios.

  Within the present framework, such sources may be interpreted as
  manifestations of primordial projective saturation.
  Two non-exclusive realizations are possible.
  First, halos exceeding $M_{\mathrm crit}(z)$ may undergo genuine direct
  collapse, producing massive compact seeds at early epochs.
  Second, a strongly constrained relational configuration may generate an
  enhanced effective convergence through the emergent contribution
  $\kappa_{\chi}$ introduced in Section~11.10 of the White Paper,
  without requiring proportionally large baryonic mass.

  In both cases, the underlying mechanism is identical:
  a resolution gap between the demanded fragmentation scale and the
  minimal projectable scale~$\ell(z)$.
  The observational signature is therefore not the presence of an exotic
  component, but the structural response of the projection under extreme
  compression.

  Discriminating between a physical direct-collapse black hole and a
  purely projective overmass configuration requires combined constraints
  from variability, spectral diagnostics, and dynamical mass estimates.
  The framework predicts that the abundance of such compact early objects
  should track the redshift evolution of the effective resolution scale
  rather than solely the mass accretion history of halos.
